\documentclass[10pt,a4paper]{amsart}
\usepackage[margin=19mm]{geometry}
\usepackage{amsmath,amssymb,mathtools}
\usepackage[scr=rsfs,cal=euler]{mathalfa}
\usepackage{xfrac}
\usepackage[unicode,hidelinks,bookmarks=false]{hyperref}
\newcommand{\ie}{i.e.\ }
\newcommand{\cf}{cf.\ }
\newcommand{\resp}{respectively\ }
\newcommand{\Subsection}{\subsection}
\providecommand{\nobreakdash}{\nobreak\hbox{-}}
\setlength{\emergencystretch}{2em}
\multlinegap0pt
\usepackage{bm}
\usepackage{bbm}
\usepackage{dsfont}
\usepackage{xspace}
\usepackage{cancel}
\usepackage{mathtools}
\usepackage{amsthm}
\makeatletter
\@ifundefined{theorem}{\newtheorem{theorem}{Theorem}}{}
\@ifundefined{lemma}{\newtheorem{lemma}[theorem]{Lemma}}{}
\makeatother
\def\N{\mathop{\mathbb N\kern 0pt}\nolimits}
\def\R{\mathop{\mathbb R\kern 0pt}\nolimits}
\def\Z{\mathop{\mathbb Z\kern 0pt}\nolimits}
\def\cQ{\mathcal Q}
\def\f{\frac}
\def\ls{\lesssim}
\def\p{\partial}
\def\ve{\varepsilon}
\newcommand{\pk}{\dot{P}_k}
\title[Long time smooth solutions of 3D cubic quasilinear wave syst]{Long time smooth solutions of 3D cubic quasilinear wave systems with small weakly decaying initial data}
\author{Mu Gao, Jun Li, Huicheng Yin}
\date{Source: arXiv:2604.17683v1 (CC BY 4.0)}
\begin{document}
\maketitle

\noindent\emph{Source and licence.} Long time smooth solutions of 3D cubic quasilinear wave systems with small weakly decaying initial data. Version arXiv:2604.17683v1 (CC BY 4.0), distributed under the Creative Commons Attribution 4.0 International licence, as recorded in the arXiv licence metadata for this exact version and shown on the abstract page https://arxiv.org/abs/2604.17683v1. Full rights notice: licenses/arxiv-2604.17683-CC-BY-4.0.txt.\par\smallskip

\noindent\textit{Editorial scope.} Counted unit: the Lemma whose source label is \texttt{lem:Energy estimate} in the article (source0.tex lines 2194--2202, source label \texttt{lem:Energy estimate}), delivered together with the article's own complete original proof of it (source0.tex lines 2205--2335, one \texttt{proof} environment of 131 lines). No mathematics is written, repaired or re-derived: statement and proof are the article's own text line for line. Required context retained uncounted: initial:data (lines 266--268); initial:data2 (lines 282--284); eq:wave2 (lines 2134--2138); eq:symmetric condition2 (lines 2152--2154). Every definition, display and label of those lines is retained unchanged. Omitted from this package and not redistributed: 1--265, 269--281, 285--2133, 2139--2151, 2155--2193, 2203--2204, 2336--3510. Nothing inside a retained excerpt is omitted or altered: every retained body is delivered exactly as the article's lines are written, so no omission receipt of any kind accompanies this package. Citations: the retained excerpts carry no citation command at all, so no bibliography entry of the article is transcribed and no reference is redistributed. Reference adaptation: none was needed and none was made; every reference inside the delivered unit resolves inside it, so no unresolved reference or citation is produced. Printed slips kept literal: no printed word, symbol, hypothesis or number of the counted statement or of its proof was altered. Layout and class: the article's own class is \texttt{article} with packages including amsmath, amssymb, amsthm, mathrsfs, simplewick, times, color, hep, leftidx, graphicx, float, enumerate, titletoc, epstopdf; the wrapper is the lane's pinned \texttt{amsart} editorial wrapper at 10pt on A4, which restates the theorem-like environments the retained text needs and adds the article's own macro definitions that the retained text needs (\texttt{\textbackslash N}, \texttt{\textbackslash R}, \texttt{\textbackslash Z}, \texttt{\textbackslash cQ}, \texttt{\textbackslash f}, \texttt{\textbackslash ls}, \texttt{\textbackslash p}, \texttt{\textbackslash pk}, \texttt{\textbackslash ve}), with extra wrapper packages (bm, bbm, dsfont, xspace, cancel, mathtools). The delivered numbering is the unit's own. Assets: no retained span contains a figure environment, a graphics-inclusion command, a PDF-page inclusion, a table or a TikZ picture; no image file is copied into this package, the package declares no asset dependency and no image file is redistributed with it. Licence: version arXiv:2604.17683v1 of this article is distributed under the Creative Commons Attribution 4.0 International licence, as recorded in the arXiv licence metadata for this exact version and shown on the abstract page https://arxiv.org/abs/2604.17683v1. A full rights notice is delivered at licenses/arxiv-2604.17683-CC-BY-4.0.txt. Characters: every retained excerpt is pure ASCII, so the delivered engine, XeLaTeX with its default Latin Modern text font, reports no missing character.
\begin{equation}\label{initial:data}
\|u_{0}\|_{H^{N+1}(\R^3)}+\|u_{1}\|_{H^N(\R^3)}\le\ve,
\end{equation}

\begin{equation}\label{initial:data2}
\|u_{0}\|_{H^{N+1}(\R^3)}+\|u_1\|_{H^N(\R^3)}+\sum_{|a|\le 5}\|(1+|x|^2)^{\mu/2}\p^a_x(\nabla u _0,u_1)\|_{L^2(\R^3)}\le\ve,
\end{equation}

\begin{equation}\label{eq:wave2}
	\begin{aligned}
		\square_{c_i} u^i =G^i(u,\p u,\p_x\p u) = \sum_{j=1}^m\sum_{\alpha,\beta=0}^3 \cQ^{\alpha\beta}_{ij} \partial_{\alpha\beta}^2 u^j+\mathcal{S}^i(u,\p u),
	\end{aligned}
\end{equation}

\begin{equation}\label{eq:symmetric condition2}
	\cQ^{\alpha\beta}_{ij} =\cQ^{\beta\alpha}_{ij} =\cQ^{\alpha\beta}_{ji}.
\end{equation}

\begin{lemma}[Energy estimate]\label{lem:Energy estimate}
	Let $\delta>0$ and integer $N \geq 2$. If $u$ is the solution of \eqref{eq:wave2}, where $\|\partial u(t)\|_{H^{N}(\mathbb{R}^3)} \leq \varepsilon_1$ for $t\ge 0$ and $\varepsilon_1>0$ is small, then one has
	\begin{equation}\label{3.4}
		\begin{aligned}
			\|\partial u(t)\|_{H^{N}(\mathbb{R}^3)} &\lesssim \|\partial u(0)\|_{H^{N}(\mathbb{R}^3)}
			+ \big(\sum_{k \geq -1} 2^{(1+\delta)k} \|P_k \partial^{\le 1} u\|_{L^2([0,t];L^\infty(\mathbb{R}^3))}\big)^2 \|\partial u\|_{L^\infty([0,t];H^{N}(\mathbb{R}^3))}.
		\end{aligned}
	\end{equation}
\end{lemma}

\begin{proof}
	Acting $\partial_x^a$ with $a\in\N_0^3$ on two sides of \eqref{eq:wave2} yields
	\begin{equation}\label{3.5}
		\begin{aligned}
			\square_{c_i}\partial_x^a u^i &= I_1^{a,i} + I_2^{a,i} +  I_3^{a,i},\\	
			I_1^{a,i} &= \sum_{j=1}^m\sum_{\alpha,\beta=0}^3 \cQ^{\alpha\beta}_{ij} \p_x^a\partial_{\alpha\beta}^2 u^j, \\
			I_2^{a,i} &= \sum_{j=1}^m\sum_{\alpha,\beta=0}^3\sum_{\substack{b+c=a, \\ |b| \leq |a|-1}}  \p_x^c\cQ^{\alpha\beta}_{ij} \p_x^b\partial_{\alpha\beta}^2 u^j, \\
			I_3^{a,i} &= \p_x^a\mathcal{S}^i(u,\p u),\\
		\end{aligned}
	\end{equation}
	where $I_2^{a,i}$ vanishes for $|a|=0$.

	Note that $\int_{\R^3}(\p_t f)^2 + c^2|\nabla f|^2 dx \approx \|\partial f\|_{L^2(\mathbb{R}^3)}^2 $.
Multiplying \eqref{3.5} by $\partial_t \partial_x^a u^i$ and integrating the resulting equality over
	$[0,t] \times \mathbb{R}^3$, it follows from the summation over the index $i$ that
	\begin{equation}\label{3.6}
		\begin{aligned}
			\|\partial_x^a\partial  u(t)\|_{L^2(\mathbb{R}^3)}^2 &\lesssim \| \partial_x^a\partial u(0)\|_{L^2(\mathbb{R}^3)}^2
+ \sum_{i=1}^{m}\big| \int_0^t \int_{\mathbb{R}^3} \partial_t \partial_x^a u^i I_1^{a,i} dx d\tau \big| \\
			&\quad +\sum_{i=1}^{m} \int_0^t \|\partial_t \partial_x^a u^i(\tau)\|_{L^2(\mathbb{R}^3)} \big(\|I_2^{a,i}\|_{L^2(\mathbb{R}^3)}+\|I_3^{a,i}\|_{L^2(\mathbb{R}^3)}\big) d\tau.
		\end{aligned}
	\end{equation}
	Firstly, we deal with the last term in the first line of \eqref{3.6}. Note that
	\begin{equation}\label{3.7}
		\begin{aligned}
			\cQ^{\alpha\beta}_{ij} \partial_t \partial_x^a u^i  \p_x^a\partial_{\alpha\beta}^2 u^j
			=&\partial_\alpha (\cQ^{\alpha\beta}_{ij} \partial_t \partial_x^a u^i  \p_x^a\partial_{\alpha} u^j)
			-\p_\alpha \cQ^{\alpha\beta}_{ij}\partial_t \partial_x^a u^i \partial_\beta \partial_x^a u^j\\
			&-\frac{1}{2} \partial_t (\cQ^{\alpha\beta}_{ij} \partial_t \partial_x^a u^i  \p_x^a\partial_{\alpha} u^j)
			+\frac{1}{2} \p_t \cQ^{\alpha\beta}_{ij}  \partial_\alpha \partial_x^a u^i \partial_\beta \partial_x^a u^j ,
		\end{aligned}
	\end{equation}
	where the Einstein summation convention for $i,j = 1,\cdots,m,\ \alpha,\beta=0,1,2,3$ and the symmetric condition \eqref{eq:symmetric condition2} are used in \eqref{3.7}. On the other hand, noting that if $u\in H^2(\R^3)\subset S'_h(\R^3)$, then it follows from Bernstein's inequality that
	\begin{equation*}
		\|u\|_{L^\infty(\R^3)}
		\ls
		\sum_{k\in\Z}2^{\f32 k}\|\pk u\|_{L^2(\R^3)}
		\ls
		\sum_{k\in\Z}2^{\f12 k}\|\pk\nabla u\|_{L^2(\R^3)}
		\ls
		\|\nabla u\|_{H^{1}(\R^3)}.
	\end{equation*}
	Hence, it holds
	\begin{equation}\label{3.8}
		\begin{aligned}
			\|\cQ^{\alpha\beta}_{ij}(t)\|_{L^\infty}
			\ls
			\|\partial u(t)\|_{L^\infty(\mathbb{R}^3)}^2 +\|u(t)\|_{L^\infty}^2
			\ls
			\|\partial u(t)\|_{H^2}^2
			\ls \ve_1^2.\\
		\end{aligned}
	\end{equation}
	It follows from \eqref{initial:data} (or \eqref{initial:data2}), \eqref{3.8} and integration over $[0,t] \times \mathbb{R}^3$
for \eqref{3.7} that
	\begin{equation*}
		\begin{aligned}
		\sum_{i=1}^{m}\big| \int_0^t \int_{\mathbb{R}^3} \partial_t \partial_x^a u^i I_1^{a,i} dx dt\big|
        &\lesssim \ve^2\|\partial_x^a\partial  u(0)\|_{L^2}^2 + \ve^2_1\|\partial_x^a\partial  u(t)\|_{L^2}^2 \\
		&\quad + \int_0^t  \|\partial_x^{\le 1}\partial^{\le 1}  u(\tau)\|_{L^\infty}^2 \|\partial_x^a\partial  u(\tau)\|_{L^2}^2 d\tau,
		\end{aligned}	
	\end{equation*}
	where the boundary term in \eqref{3.7} vanishes due to $\cQ^{\alpha\beta}_{ij}\partial_t \partial_x^a u^i \partial_\beta \partial_x^a u^j  \in L^1(\mathbb{R}^3)$.\\
	Next, we treat $\|I_2^{a,i}\|_{L^2(\mathbb{R}^3)}$ with $|a| \geq 1$ in \eqref{3.5}.
It is easy to find that
	\begin{equation*}
		\begin{aligned}
			\|I_2^{a,i}\|_{L^2}
			\lesssim& \big\| \|P_k I_2^{a,i}\|_{L^2}\big\|_{\ell_k^2} \\
			\lesssim& \sum_{j=1}^m\sum_{\alpha,\beta=0}^3\sum_{\substack{b+c=a, \\ |b| \leq |a|-1}} \big\| \sum_{(k_1,k_2) \in \mathcal{X}_k} \| P_k(P_{k_1} \partial_{\alpha\beta}^2 \partial_x^b u^j P_{k_2} \p_x^c\cQ^{\alpha\beta}_{ij} )\|_{L^2} \big\|_{\ell_k^2}.
		\end{aligned}
	\end{equation*}
Note that $\mathcal{Q}^{00}_{ij}=0$ and $|c|\ge 1$. When $\max\{k_1, k_2\} = k_1$, one then has that by Bernstein's inequality
	\begin{equation*}
		\begin{aligned}
			&\|P_k(P_{k_1} \partial_{\alpha\beta}^2 \partial_x^b u^j P_{k_2} \p_x^c\cQ^{\alpha\beta}_{ij} )\|_{L^2}\\
			&\lesssim \|P_{k_1} \partial\partial_x \partial_x^b u\|_{L^2} \|P_{k_2} \p_x^c\cQ^{\alpha\beta}_{ij} \|_{L^\infty} \\
			&\lesssim 2^{k_1(|b|+1)+k_2(|c|-1)} \|P_{k_1} \partial u\|_{L^2} \|P_{k_2}\p_x \cQ^{\alpha\beta}_{ij}\|_{L^\infty}  \\
			&\lesssim 2^{k_1|a|} \|P_{k_1} \partial u\|_{L^2} \|P_{k_2} \p_x \cQ^{\alpha\beta}_{ij}\|_{L^\infty}.
		\end{aligned}	
	\end{equation*}
	Therefore,
	\begin{equation*}
		\begin{aligned}
			\big\| \sum_{\substack{(k_1,k_2) \in \mathcal{X}_k, \\ \max\{k_1,k_2\}=k_1}} \| P_k(P_{k_1} \partial_{\alpha\beta}^2 \partial_x^b u^j P_{k_2} \p_x^c\cQ^{\alpha\beta}_{ij} )\|_{L^2}\big\|_{\ell_k^2}
			\lesssim \|\partial u\|_{H^{|a|}}  \| \cQ^{\alpha\beta}_{ij} \|_{B^{1+\delta}_{\infty,1}}.
		\end{aligned}
	\end{equation*}
	When $\max\{k_1, k_2\} = k_2$, we can analogously obtain
	\begin{equation}\label{3.11}
	\begin{aligned}
	\big\| \sum_{\substack{(k_1,k_2) \in \mathcal{X}_k, \\ \max\{k_1,k_2\}=k_2}} \|P_k(P_{k_1} \partial_{\alpha\beta}^2 \partial_x^b u^j P_{k_2} \p_x^c\cQ^{\alpha\beta}_{ij} )\|_{L^2}\big\|_{\ell_k^2}
	\lesssim \|\p_x\cQ^{\alpha\beta}_{ij} \|_{H^{|a|-1}} \| \partial u\|_{B^{1+\delta}_{\infty,1}} .
	\end{aligned}	
	\end{equation}
Note that both $\dot{B}^s_{p,r}(\R^3)\cap L^\infty(\R^3)$ and $B^s_{p,r}(\R^3)\cap L^\infty(\R^3)$ are algebras for $s > 0$, $p,r\in[1,\infty]$.
Then for $1\le |a|\le N$, it holds
\begin{equation*}
\begin{split}
\| \cQ^{\alpha\beta}_{ij} \|_{B^{1+\delta}_{\infty,1}}&\ls \| \p^{\le 1} u \|_{B^{1+\delta}_{\infty,1}}\| \p^{\le 1} u \|_{L^{\infty}}\ls \| \p^{\le 1} u \|_{B^{1+\delta}_{\infty,1}}^2,\\
\|\p_x\cQ^{\alpha\beta}_{ij} \|_{H^{|a|-1}}&\ls\|\p_x\cQ^{\alpha\beta}_{ij} \|_{L^{2}} +\|\cQ^{\alpha\beta}_{ij} \|_{\dot{H}^{|a|}}\\
&\ls\|\partial_x^{\le 1}\partial  u\|_{L^2}\|\partial^{\le 1}  u\|_{L^\infty}
+\|\partial^{\le 1}  u\|_{\dot{H}^{|a|}} \|\partial^{\le 1}  u\|_{L^\infty}\\
&\ls\|\partial u\|_{{H}^{|a|}} \|\partial^{\le 1}  u\|_{L^\infty}.
\end{split}
\end{equation*}
Hence,
\begin{equation*}
\sum_{i=1}^{m}\|I_2^{a,i}\|_{L^2}\ls \|\partial u\|_{H^{|a|}} \|\partial^{\le1} u\|_{B^{1+\delta}_{\infty,1}}^2.
\end{equation*}
Finally, we handle $\|I_3^{a,i}\|_{L^2(\mathbb{R}^3)}$ in \eqref{3.5}. For $|a|= 0$, one has
\begin{equation*}
\sum_{i=1}^{m}\|I_3^{0,i}\|_{L^2}\ls \|\p u\|_{L^2}\|\p^{\le 1} u\|_{L^\infty}^2.
\end{equation*}
For $1\le|a|\le N$, it holds
\begin{equation}\label{3.15}
\begin{split}
\sum_{i=1}^{m}\|I_3^{a,i}\|_{L^2}
&\ls \|\partial u\|_{\dot{H}^{|a|}} \|\partial^{\le1} u\|_{L^\infty}^2 + \|\partial^{\le 1} u\|_{\dot{H}^{|a|}} \|\partial^{\le1} u\|_{L^\infty}\|\partial u\|_{L^\infty}\\
&\ls 	\|\partial u\|_{{H}^{|a|}} \|\partial^{\le 1}  u\|_{L^\infty}^2.
\end{split}	
\end{equation}
Collecting \eqref{3.6}--\eqref{3.15} together with the smallness of  $\ve$ and $\ve_1$ yields
\begin{equation}\label{4.29}
\begin{aligned}
\|\partial  u(t)\|_{H^{|a|}(\mathbb{R}^3)}^2 &\lesssim \|\partial  u(0)\|_{H^{|a|}(\mathbb{R}^3)}^2
+\int_0^t   \|\partial^{\le1} u(\tau)\|_{B^{1+\delta}_{\infty,1}(\R^3)}^2   \|\partial  u(\tau )\|_{H^{|a|}(\mathbb{R}^3)}^2 d\tau.
\end{aligned}
\end{equation}
Then \eqref{3.4} is derived  from \eqref{4.29}, Gronwall's inequality and H\"older's inequality.
\end{proof}

\end{document}
